% II. Проектиране и програмна реализация.
\section{Проектиране и програмна реализация}
\label{sec:design}

\subsection{Обща структура и модел на изпълнение}

Сървърният процес чете показатели от операционната система, оформя ги в стойностна снимка,
превръща снимката в съобщение и го изпраща по избрания от клиента канал. Страницата в браузъра
приема съобщението, разчита го и обновява само засегнатите части от документа
\figref{fig:architecture}. Сървърът е разделен на шест единици за превод: \code{net},
\code{http}, \code{websocket}, \code{metrics}, \code{codec} и \code{server}. Разделянето следва
посоката на данните, а не слоевете на протокола.

\begin{figure}[H]
\centering
\begin{tikzpicture}[x=1cm, y=1cm]
  \node[block, text width=3.4cm] (col)  at (0, 0)    {\code{metrics::Collector}\\извадка на интервал};
  \node[block, text width=3.4cm] (snap) at (4.4, 0)  {\code{metrics::Snapshot}\\една снимка};
  \node[accent, text width=2.9cm] (json) at (8.8, 1.0) {\code{codec::to\_json}};
  \node[accent, text width=2.9cm] (xml)  at (8.8, -1.0) {\code{codec::to\_xml}};
  \node[block, text width=3.4cm, minimum height=1.6cm] (loop) at (13.2, 0)
    {\code{pulse::Server}\\цикъл върху \code{select},\\маршрути, обратно налягане};
  \node[note, text width=3.6cm] (client) at (18.0, 0)
    {\textbf{Страница в браузъра}\\[3pt]
     разчитане според формата\\
     свързване на данни\\
     текстови възли, SVG};

  \draw[line] (col)  -- (snap);
  \draw[line] (snap) -- (json);
  \draw[line] (snap) -- (xml);
  \draw[line] (json) -- (loop);
  \draw[line] (xml)  -- (loop);

  \draw[line] ([yshift=5mm]loop.east) -- node[above, font=\scriptsize] {HTTP/1.1}
    ([yshift=5mm]client.west);
  \draw[line] ([yshift=0mm]loop.east) -- node[above, font=\scriptsize] {WebSocket}
    ([yshift=0mm]client.west);
  \draw[line] ([yshift=-5mm]loop.east) -- node[above, font=\scriptsize] {събития}
    ([yshift=-5mm]client.west);
  \draw[line] ([yshift=-10mm]client.west) -- node[below, font=\scriptsize] {команди}
    ([yshift=-10mm]loop.east);
\end{tikzpicture}
\caption{Обща структура на системата и посока на данните}
\label{fig:architecture}
\end{figure}

Носещото проектно решение е, че \term{сериализаторът е отделен от източника на показатели}.
Един и същ поток се излъчва в разширяемия език за разметка~\cite{xml10} или в нотацията за
обекти~\cite{rfc8259}, без нито един ред от източника или от цикъла да се променя. Без това
разделяне сравнението между двата формата щеше да сравнява две реализации, а не два формата.

Сървърът е \term{еднонишков}. Един цикъл наблюдава слушащия сокет и всички връзки с
\code{select}, обслужва готовите и на всеки интервал взема нова извадка и я разпраща. Няма
работни нишки и няма ключалки, защото всяко споделено състояние се пипа от точно една нишка.
Цената е заявена предварително: бавен обработчик би спрял всички клиенти, което е приемливо,
защото нито един обработчик не прави друго освен да оформи низ от стойности, които вече са в
паметта. Всяка връзка носи собствени буфери и режим: HTTP, надградена по WebSocket или отворен
поток от събития. Режимът се сменя в средата на буфера.

\subsection{Договор между сървъра и клиента}

От сървъра идват три вида съобщения: стойностна снимка, потвърждение на команда и съобщение
за грешка. Снимката носи пореден номер, времеви отпечатък, натоварване на процесора, заета
памет, броячи на заявки и на връзки, три квантила и хистограма от девет коша. От клиента идват
три команди: \code{format} със стойност \code{json} или \code{xml}, \code{subscribe} и
\code{unsubscribe}, и \code{interval} между 50 и 5000 милисекунди. Командите вървят само по
WebSocket, защото другите два транспорта са еднопосочни. Отпечатъкът е нужен само за
измерването: клиентът го изважда от собствения си часовник и получава възрастта на стойността
при пристигане.

Едно и също съдържание се раздава по три начина: \code{/ws} отговаря с 101 или 426,
\code{/api/stream} с парчетно кодиран поток от събития, а \code{/api/metrics.json} и
\code{/api/metrics.xml} с последната снимка за периодично запитване. Отговорът 426 не е грешка
в клиента, а заявка, която сървърът не може да изпълни както е написана; стандартът има точно
това състояние за случая~\cite{rfc9110}.

\subsection{Свързване на данни на клиента}

Свързването е явен списък от тройки: път до стойност в снимката, възел от документа и
функция, която записва стойността в този възел. При пристигане на снимка се обхождат само
тройките с променена стойност, защото при няколко обновявания в секунда пълното преизчисляване
на изгледа прави повече работа, отколкото има променени стойности. Диаграмите се изграждат като
структурна графика~\cite{svg2}: всяка точка от реда е възел в документа и приема обработчик на
събитие и стилов преход, а разположението е зададено с CSS Grid~\cite{css_grid1}. Обработката
на събития е съсредоточена чрез \term{делегиране} върху общия родителски възел, защото възлите
на диаграмата се създават и унищожават при всяко обновяване.

Проектът приема четири ограничения, заявени преди да са измерени резултати: състоянието се пази
в паметта на процеса и историята се губи при рестарт; поддържа се един източник на показатели;
няма удостоверяване, което е приемливо само защото системата се изпълнява локално; интервалът
на публикуване е общ за всички клиенти.
